%% Evaluation protocol — corpus details cross-ref sec:data
\subsection{Corpus and indexing grid}
As detailed in~\cref{sec:data}, retrieval experiments use a fixed IAEA \emph{core} pool of 13 documents
and a same-domain \emph{non-core} (supplementary) pool of up to 46 manuals (59 documents in total).
We vary how many non-core documents are admitted alongside the core set, thereby scanning different
core--non-core ratios rather than treating non-core texts as unstructured noise.
Seventeen Markdown index configurations span chunk sizes $\{800,1000,1200,1500\}$ and overlaps
$\{50,80,100,120,150\}$ (valid pairs only).
The embedding model is fixed (\texttt{nomic-embed-text}); only the active retrieval pool and chunking
configuration change across runs.

\subsection{Experiment A: corpus-scope extrapolation}
Four fixed queries---Chinese/English core terms and phrases (e.g., ``neutron'' /
``neutron measurement current'')---are issued while randomly attaching non-core documents in steps
of five (multiple trials), so that the effective core:non-core composition changes in a controlled
manner.
We report top-$k$ cosine scores, core (IAEA) purity@10 (fraction of top-10 hits from the core pool),
and the rank of the first core hit.

\subsection{Experiment B: self-retrieval}
IAEA chunks are sampled as probes.
Queries are built as first sentence, title, random window, or domain-term sentence.
Relevance is scored under \emph{strict} (same \texttt{node\_id}) and \emph{relaxed\_doc} (same
\texttt{doc\_id}) criteria.
Metrics: Recall@$\{1,5,10\}$ and MRR, including decay under corpus expansion and the strict--relaxed
gap induced by chunk overlap.
